%% The line mask, reported in two layers: the frequency coincidence, which is
%% a velocity criterion and nothing else, and the one species we do not read
%% as emission, which is an argument and is set out as one.
%% Owner: r16-mask.  Carries \label{app:linemask} (cited by 04_method.tex and
%% app_C_spatialnorm.tex) and \label{tab:masklines}.
%%
%% ★ THIS APPENDIX IS THE PAPER'S ONLY DEFINITION OF THE ATTRIBUTION RULE.
%% It owns the species list, the half-width and the FRAME.  The rule is the
%% star's own rest frame; the local standard of rest is evaluated as a check
%% and changes no disposition (masklsr_v520.py, L10).  Sec. 5.3 states the
%% rule in one clause and Appendix D points here; nothing else in the paper
%% may name a frame of its own, or there are two rules again.
\section{The line mask}
\label{app:linemask}

A crossing that falls on a known molecular transition has an astrophysical
explanation to hand, and the mask exists to say which crossings those are. It
reports the frequency coincidence and the judgement about the coinciding
species separately, so the reproducible part can be recomputed and the
arguable part argued with.

\emph{The list, and when it was made.} The mask holds every transition of the
\FrNSpec{} species these spectral setups cover -- the three CO isotopologues,
HCN, HCO$^{+}$, CS, CN, SiO, H$_2$CO, SO, [C\,\textsc{i}] and hydrogen
recombination -- that falls inside one of the \McNIsland{} searched frequency
intervals with an upper-state energy below \McEuMax\,K. That rule is the whole
of it, and it admits \FrNTrans{} transitions (Table~\ref{tab:masklines}).
No strength cut is applied across species: any threshold admitting HCN's
strong lines would discard the brightest lines in these fields. Within a
species the main hyperfine components are kept and the satellites
dropped. The list is a single Splatalogue query of CDMS and
JPL \citep{Pickett1998,CDMS2005} made on \MkListDate{}, after the disposition
criteria were fixed on \MkCritDate, and frozen before any crossing was
compared with it.

\emph{The species are the ones the observers went to the telescope for, and
that is a limitation.} The spectral setups are tuned to them. \TnCenPct{} per
cent of the \NWinA{} Class~A windows are centred within
$\pm\MaskHalfKms$\,km\,s$^{-1}$ of a transition of the list and \TnWinPct{}
per cent contain one somewhere. Displace the same windows bodily by
\TnShiftLoGHz{} to \TnShiftHiGHz\,GHz and only
\TnCtrlCenLo--\TnCtrlCenHi{} per cent remain centred on one. The list is
therefore set by the science goals of the
programmes that took the data and not by anything in this search, which is
why it works as well as it does on these fields and also where it is blind: a
setup tuned to a line the list does not hold is screened only where that line
and a masked transition happen to coincide in frequency.

\emph{Attribution is decided in the star's own rest frame, and the sign is
not the usual one.} A crossing's frequency is carried from the topocentric
value to the barycentre using its own block's pointing and mid-time, and then
to the stellar rest frame using the star's systemic velocity. That frame, and
no other, is where a coincidence is judged.
The barycentric term alone spans \FrBarySwing\,km\,s$^{-1}$ across the
survey, so a tolerance applied in the observed frame would spend most of its
budget on the Earth's motion. The \FrEpochNEpoch{} \FrEpochStar{} epochs show
the transform working in the only way that can be checked:
\FrEpochTopoKms\,km\,s$^{-1}$ apart as observed, they agree to
\FrEpochStelKms\,km\,s$^{-1}$ in the star's frame.
Throughout, $\Delta v_\star = c\,(\nu_\star-\nu_0)/\nu_0$, so a positive
offset means the crossing lies \emph{above} the transition in frequency --
the opposite sign to the velocity convention of a spectral line.

\emph{Every crossing is read in the local standard of rest as well, where
foreground gas is at rest, and what that second reading establishes is which
coincidences are foreground.} Of the \LsNLos{} stars carrying crossings, the composite
CO($1\to0$) survey of \citet{DameThaddeus2022}, which contains
\citet{Dame2001}, covers \LsNLosCov{} and holds emission toward \LsNLosDet{}
of them, between \LsCoVLo{} and \LsCoVHi\,km\,s$^{-1}$. Toward every other
star, including the \LsNLosOut{} outside the CfA footprint, the all-sky
Planck CO map \citep{PlanckCO2016} stays below
\LsPlanckMax\,K\,km\,s$^{-1}$. Toward
\LsCloudStars{} the CO spectrum is non-zero at the crossing's own LSR
velocity, within \LsCloudDvMax\,km\,s$^{-1}$ of its own peak and at
\LsCloudTMin\,K or brighter against a \LsCoRms\,K channel. So \LsNCloud{} of
the \MkNCoinc{} coincidences are foreground gas and not circumstellar.
Two are worth naming. \BkAlmaName{} is a serendipitous field star toward the
Lupus clouds, at \BkAlmaDistPc\,pc with the clouds far behind it. Its
CO($2\to1$) crossing lies \LsLupDv\,km\,s$^{-1}$ from a
\LsLupWco\,K\,km\,s$^{-1}$ cloud line at \LsLupVco\,km\,s$^{-1}$, so the
field is a molecular sightline and not a disc. The velocity says so and the
field
does not: its brightest control position stands above the star in all
\LsLupNRingLeadWord{} of its windows, but so it does in \LsRingLeadA{} of the
\NWinA{} Class~A windows of the survey, so that comparison
distinguishes nothing. Its other
\LsLupNFar{} crossings lie at least \LsLupFarMin\,km\,s$^{-1}$ from every
transition of the list, and the mask disposes of none of them. HD~48370 is
the case where the two readings part: its coincidence lies
\LsHdDvStar\,km\,s$^{-1}$ from the star's own CO($2\to1$), while in the local
standard of rest it sits at \LsHdVfeat\,km\,s$^{-1}$, on a catalogued cloud
line peaking at \LsHdVco. The offset follows the foreground gas and not the
star, so that window's extended emission is foreground cloud and not a
resolved disc.

\emph{A coincidence is an attribution, and for most of them it is not the
only evidence.} \LnAttrSupp{} of the \MkAttr{} attributions carry support
beyond the frequency match: catalogued carbon monoxide on the sightline here,
a second transition of the same species, or the crossing's own recurrence.
The other \LnAttrCoincWord, toward \LnCoincStars, rest on the coincidence
alone -- no catalogued gas on either sightline, no second transition and no
repeat epoch that recovers them -- so they are consistent with molecular
emission and are not identified as it.

\begin{table}
\centering
\footnotesize
\caption{The line mask, by species: the number of transitions $N$ it holds,
the frequency range they span inside the searched intervals and their
upper-state energies. Hydrogen recombination lines carry no tabulated
$E_{\rm u}$.}
\label{tab:masklines}
%% GENERATED by maskso_v414.py -- do not hand-edit.
\begin{tabular}{@{}lrr@{\,}c@{\,}rr@{}}
\hline
Species & $N$ & \multicolumn{3}{c}{$\nu_0$ (GHz)} & $E_{\rm u}$ (K) \\
\hline
H & 5 & 92.03 & -- & 231.90 & -- \\
SO & 15 & 100.03 & -- & 682.68 & 16--143 \\
H$_2$CO & 13 & 101.33 & -- & 491.97 & 10--106 \\
C$^{18}$O & 3 & 109.78 & -- & 329.33 & 5--32 \\
CN & 50 & 113.12 & -- & 680.30 & 5--114 \\
CO & 3 & 115.27 & -- & 345.80 & 6--33 \\
CS & 3 & 195.95 & -- & 342.88 & 24--66 \\
SiO & 2 & 217.10 & -- & 347.33 & 31--75 \\
$^{13}$CO & 8 & 220.40 & -- & 330.59 & 16--32 \\
HCN & 7 & 354.50 & -- & 354.51 & 43--43 \\
HCO$^{+}$ & 1 & 356.73 & -- & 356.73 & 43 \\
{[C\,\textsc{i}]} & 1 & 492.16 & -- & 492.16 & 24 \\
\hline
\end{tabular}%

\end{table}

\emph{The coincidences, and the one species we do not read as emission.} Of
the \NCross{} crossings, \MkNCoinc{} fall within
$\pm\MaskHalfKms$\,km\,s$^{-1}$ of a transition of the list in their own
star's frame; that count is the first layer and all the velocity criterion
decides. Coincidences with sulphur monoxide account for \MkNSODecide{} of them:
\MkSOStars, at \MkSODvList\,km\,s$^{-1}$ from transitions of upper-state
energy \MkSOEuLo--\MkSOEuHi\,K, each inside a Band~7 window that also covers
CO($3\to2$). Two things argue against reading those as emission.
Sulphur monoxide has not been reported in a debris disc, where the
second-generation gas is CO, C and O. And toward \FrCpStar{} the CO($3\to2$)
line falls in the same window and is absent to \FrCpCoLimMJy\,mJy in a
\FrCpCoChanKms\,km\,s$^{-1}$ channel, while the crossing's fitted amplitude is
$\RcCpFlux\pm\RcCpFluxErr$\,mJy: an SO line there would have to be at
least as bright as an undetected CO($3\to2$) line in the same data. That
margin is only \RqAmpSigmaCp{} of the fitted error, so it is the species and
not the margin that carries the weight. Both arguments are about the species
and
not about any one crossing, so all \MkNSODecide{} are treated alike:
the ledger prints the offset beside each with a dagger and
disposes of none of them, giving \MkAttr{} attributed against \MkUnattr{}
unattributed, of which \MkRankLeadWord{} outrank every control position in
their own windows.
Reading them as attributions instead gives \MkAttrAlt{} and \MkUnattrAlt,
and \MkRankLeadAlt{} such crossing; not one of the \MkNSODecide{} is present
again in any block covering it, and none is confirmed, so the choice moves
counts and no disposition. The ledger forms the exclusion of
Eq.~\ref{eq:excl} for each of them, at
\MkAttrExclLo--\MkAttrExclHi$\sigma$, but that significance is built on the
discovery amplitude; read at the trigger amplitude instead, a crossing whose
repeat leaves the criterion without power is reported in
\S\ref{sec:etacrv} as untested rather than excluded, and that applies to
these as to any other. The reserved hold-out is treated identically, leaving
\HoNAttr{}
attributed and \HoNUnattr{} unattributed there (\S\ref{sec:holdoutsearch}).

\emph{The half-width, and what it costs.} A one-channel test is too tight for
circumstellar gas, whose emission is shifted by many channels: the width
follows from Keplerian motion where these molecules survive,
$\lesssim$13\,km\,s$^{-1}$ beyond 10\,au, and is
$\pm\MaskHalfKms$\,km\,s$^{-1}$. At that width the mask removes
\MkMaskAGHz\,GHz of the \CvUnionA\,GHz Class~A union, \MkMaskAPct{} per cent
of it. Over both classes together it removes \FrMaskLostGHz\,GHz of
\FrMaskUnionGHz\,GHz, \FrMaskPct{} per cent, so the cost falls on the primary
experiment, because its windows are the line-tuned ones.
Attributing on either frame instead would mask the union of two tubes
separated by the star's LSR velocity, \LsMaskTwoGHz\,GHz or
\LsMaskTwoPct{} per cent of the Class~A union -- \LsMaskAddPct{} per cent
more bandwidth for no further attribution, which is the measurement behind
deciding in one frame. Between $\pm\FrLadderLo$ and
$\pm\FrLadderHi$\,km\,s$^{-1}$ the mask attributes \FrAttrWTwenty{} to
\FrAttrWHundred{} of the crossings. The \FrNRankPass{} that outrank all
\NCtrl{} control positions of their own windows are the same two at every
width.

\emph{The labelled crossings are tested for recurrence too.} \MkAttrTested{}
of the \MkAttr{} attributed crossings have a later block covering the
predicted cell, and the \MkAttrRecur{} that are present again, toward
\MkAttrRecurStars, are the positive control of that stage. For the
\MkAttrExcl{} not recovered a repeat window's own largest statistic bounds
the predicted cell from above and the ledger prints the row as a
bound; \MkAttrUntest{} have no covering block at all
(\MkAttrUntestNames).
